\subsection{Sequential Containers}

\subsubsection{Vectors}

\textbf{Definiton:}\\
\texttt{vector<T> v;} 

\textbf{Element access}

\begin{tabularx}{\linewidth}{lX}
\texttt{operator[]}	& \\
\texttt{v.back()} & Last element \\
\texttt{v.front()} & First element \\
\end{tabularx}


\textbf{Modifiers}

\begin{tabularx}{\linewidth}{lX}
\texttt{v.push\_back(5)} & Add \texttt{5} to the vector. Simillar to push in stack data structure.\\
\texttt{v.insert()} & Inserts elements / range in \texttt{v}.\\
& Eg: Append \texttt{u} to \texttt{v}.\\
& \texttt{v.insert(v.end(), u.begin(), u.end());}\\
\texttt{v.emplace} & \texttt{v.emplace\{iterator pos, Args\}}. New element is constructed in place using \texttt{Args} and inserted at pos.\\
\texttt{v.emplace\_back} & \texttt{v.emplace\_back\{Args\}}. New element is constructed in place using \texttt{Args} and inserted at the end.\\
& \textcolor{red}{IMPORTANT: Note the differences between pushback, insert, emplace and emplaceback}.\\
\texttt{v.pop\_back()} & Remove element on top of the stack.\\
\texttt{v.erase()} & Remove element / range from \texttt{v}.\\
\texttt{v.clear()} & Clears all the elements in \texttt{v}.\\
\end{tabularx}

\textbf{Capacity}

\begin{tabularx}{\linewidth}{lX}
\texttt{v.size()} & Returns the size of the vector.\\
\texttt{v.capacity()} & Current memory allocated to \texttt{v} in terms of number of elements.\\
\texttt{v.max\_size()} & Maximum memory available for \texttt{v}.\\
\texttt{v.reserve()} & Reserve a minimum size for v. Does not affect v.size().\\
\texttt{v.shrink\_to\_fit()} & Shrink the capacity to fit the size.\\
\texttt{v.empty()} & Test if \texttt{v} is empty.\\
\end{tabularx}

\vspace*{8pt}
\texttt{tie} Can be used to unpack a vector.\\
\textit{Example:} \texttt{tie(a, b, c) = v;}\\


\vfill \null
\columnbreak


\subsubsection{Queue}
Eg: \texttt{queue<int> q;}\\

\begin{tabularx}{\linewidth}{lX}
\texttt{q.push(5)} & Adds element to the end.\\
\texttt{q.pop()} & Removes first element.\\
\texttt{q.front()} & Returns first element.\\
\texttt{q.empty()} & True if q is empty.\\
\end{tabularx}

\subsubsection{Dequeue}
Eg: \texttt{dequeue<int> d;}\\

\begin{tabularx}{\linewidth}{lX}
\texttt{d.push\_back(5)} & Add to the top.\\
\texttt{d.push\_front(2)} & Add to the bottom.\\
\texttt{d.pop\_back()} & Remove from the top.\\
\texttt{d.pop\_front()} & Remove from the bottom.\\
\end{tabularx}

